% 第 23 章　变化的磁场会制造电场
\chapter{变化的磁场会制造电场}

\step{反常识开场}

家里停电的夜晚，你也许用过一种应急手电筒：它里面没有电池，握在手里使劲摇晃几十下，或者扳动一个手柄转上几十圈，灯就亮了，而且能亮好几分钟。

请先别急着翻后面的解释，认真体会一下这件事有多奇怪。你学过的电路知识说：让小灯泡发光，需要电源提供电压，驱动电荷在闭合回路里流动（第 21 章）。电池靠化学过程维持电势差，可是这把手电筒里没有电池。你也没有往里面灌任何“电”，你只是摇了摇它——也就是说，你付出的只是力气，是机械运动。机械运动是怎么变成电流的？

更奇怪的还在手感上。如果你拆开过这种手电筒（或者玩过手摇充电器），你会发现摇动时手柄明显“发涩”，有一种看不见的阻力在跟你对着干，而且灯接得越亮，摇起来越费力。阻力是从哪里来的？手柄只是带动一块磁铁在线圈附近转动而已，磁铁和线圈之间明明什么都没有接触。

类似的事情还有：电磁炉的玻璃面板摸上去只是温温的，放在上面的铁锅却能烧开水；手机放在无线充电板上，两者之间没有一根导线相连，电却“隔空”进了手机。这几件事背后是同一条物理规律，它是整个电气时代真正的地基：没有它，就没有发电机，没有电网，你家里插座里的电也无从谈起。

\step{先猜一猜}

在继续读之前，请把你的猜测写下来——写得越具体越好，后面可以对照检查。

关于手电筒，常见的猜测有这么几种。猜测一：里面其实藏了一节小电池或者一个发条，摇动只是幌子。猜测二：摇动时的摩擦产生了静电，就像脱毛衣时的火花。猜测三：磁铁在线圈里移动时，“吸”出了线圈里的电荷。猜测四：运动的磁铁以某种方式推动了线圈里的电荷。

再猜一个相关问题。找一根铜管或铝管（铜和铝都不被磁铁吸引，这一点可以先用冰箱贴验证），让一块小磁铁从管子里竖直掉下去，再让一块同样大小的普通铁块掉下去做对比。你预测：磁铁下落会和铁块一样快，还是更快，还是更慢？大多数人凭直觉会选“一样快”——铜铝又不吸磁铁，管子里除了空气什么也没有，还能有什么差别？

最后猜一个定量的问题：如果摇动得快一倍，手电筒发出的电会多一点、多一倍、多好几倍，还是不变？

\step{动手实验}

\begin{experiment}[磁铁掉进金属管的慢动作]
\textbf{材料}：一块钕磁铁（旧电脑硬盘里拆出来的两块强磁就很好，网购直径一厘米左右的小圆柱磁铁也只要几块钱）；铝箔纸，卷成直径约两厘米、长约十五厘米的管子，卷三四层让它足够挺括；一块大小差不多的非磁性物体，比如一节同样的铝箔卷成的实心卷或者塑料块。注意：钕磁铁夹手很疼，请远离手机、银行卡和手表。

\textbf{步骤}：先用磁铁吸一吸铝箔，确认它们之间没有任何吸引。把铝箔管竖直拿在手里，先把非磁性物体从管口放进去——它“啪”的一声直接落到底。现在把钕磁铁从同样的高度放进去。

\textbf{你会看到}：磁铁仿佛换了一个人，慢悠悠地、几乎是优雅地往下飘，花上好几秒才落到管底。换一根更粗的管子或者用铜管，效果类似；把管子横过来倾斜一个角度，磁铁也是缓慢地滚下去。

\textbf{请记录}：磁铁和铝明明互不吸引，是什么在“托着”磁铁？磁铁被托住而不加速，它少掉的那些动能（或者说重力势能）到哪里去了？记住你的疑问，本章结束时你要能完整回答它。
\end{experiment}

\step{旧模型遇到麻烦}

现在把手头已有的工具摆出来，看看它们够不够用。到上一章为止（第 22 章），你已经知道三件事：电荷周围存在电场，电场对电荷施力；电流和运动电荷周围存在磁场；磁场对运动电荷施加洛伦兹力，方向既垂直于电荷的运动方向、又垂直于磁场方向，所以磁场通常不直接对电荷做功。

先检查前面的猜测。猜测一（藏电池）可以被实验排除：拆开看，里面只有一块磁铁、一个线圈、几个小元件和一颗小电容。猜测二（摩擦静电）也可以排除：摩擦起电能产生几千伏的电压，但只能放出头发丝粗细、转瞬即逝的火花，绝不可能持续点亮一颗灯好几分钟——静电给不出持续的能量流。猜测三（磁铁“吸出”电荷）说不过去：磁铁吸铁、钴、镍，不吸铜，而手电筒的线圈恰恰是铜做的。

剩下猜测四：磁铁的运动推动了线圈里的电荷。这个方向对了，但请往下追问一层：推动的机制是什么？线圈里的自由电荷本来安安静静待在铜原子之间，速度几乎为零。而洛伦兹力的大小正比于电荷的运动速度——速度为零，磁力就是零。一个对静止电荷完全无能为力的磁场，凭什么让静止电荷动起来、排成队定向流动？

有人会说：可是磁铁在动啊，是磁铁的“动”带动了电荷。那好，请看法拉第在 1831 年做的那个决定性实验。他用一个软铁圆环，一边绕上线圈 A 接电池和开关，另一边绕上线圈 B 接电流计，两个线圈互不接触，唯一的联系是穿在同一个铁环里。合上开关的瞬间，线圈 A 里电流从无到有，铁环里的磁场从无到有——就在这一瞬间，线圈 B 的电流计指针摆了一下；开关保持合上、磁场稳定不变，指针回到零；断开开关、磁场消失的瞬间，指针又摆了一下，方向相反。

请特别注意：整个实验里没有任何东西移动。线圈 B 是静止的，它里面的电荷也是静止的，唯一发生的事情是穿过它的磁场\emph{变了}。洛伦兹力解释不了它——静止电荷不受磁力。旧的三件工具全部用不上。我们遇到了一个真正的裂缝：\textbf{静止的电荷，在变化的磁场旁边，开始定向移动。一定有一种我们还没命名的力在推动它。}

\step{发明新概念}

法拉第本人没有受过多少数学训练，他用一种非常“可视化”的方式思考：他想象磁铁的周围空间里布满了“磁感线”，铁屑在磁铁周围排成的纹路就是这种线的直接证据。然后他提出了一个改变世界的观察角度：不要去管磁铁动不动、线圈动不动，只问一件事——\textbf{穿过一个闭合回路的磁感线“总量”有没有变化？}

这个“总量”后来被命名为\textbf{磁通量}。你可以这样建立直觉：把回路想象成一个圆形的渔网，把磁场想象成水流，磁通量就是单位时间里穿过网眼的水量——水量越大、网张得越正，穿过去的就越多。这个比喻像的地方在于：它正确地告诉你，通量取决于三个因素——磁场有多强、回路面积有多大、回路的朝向有多“正”。它不像的地方在于：磁场里并没有任何东西真的在流动。“磁通量”不是某种流体的流量，它是对空间中磁场分布的一笔总账，一个纯粹的数学记账方式。没有东西流过你的渔网，账却实实在在地记在那里。

法拉第的结论是：只要这笔账变了——不管你用什么方法让它变，磁铁靠近、线圈移动、回路转动、回路面积伸缩，或者像软铁环实验那样让磁场本身增强减弱——回路里就会出现推动电荷流动的力量。电荷被推动了，说明回路中出现了电场。这个电场就是我们为裂缝找到的新概念：\textbf{感生电场}，也叫\textbf{涡旋电场}。它是变化的磁场在空间中直接制造出来的，不需要任何电荷当“源头”。

请注意它和第 19、20 章里静电场的深刻区别，这个区别值得列出来慢慢消化：

\begin{itemize}
\item \textbf{来源不同}：静电场由电荷产生；涡旋电场由变化的磁场产生，空间中可以一个电荷都没有。
\item \textbf{形状不同}：静电场的电场线从正电荷出发、到负电荷终止，有头有尾；涡旋电场的电场线没有头尾，是绕着磁场变化区域一圈一圈的闭合圆环，像水面上的漩涡。
\item \textbf{绕一圈的账不同}：在静电场里，让一个电荷沿任意闭合路径走一圈，电场做的总功严格等于零——正因为这样，我们才能给每一点定义一个确定的电势（第 20 章）。而涡旋电场里，电荷绕闭合回路一圈，电场持续地推它，总功\emph{不为零}。这个“绕一圈的总推动”就是我们说的\textbf{感应电动势}。
\end{itemize}

最后一条有个重要推论，请记牢：感应电动势是\emph{沿着整个回路}定义的量，不是某两点之间的电势差。在涡旋电场里，“这一点的电势是多少”这个问题本身就没有唯一答案——从同一点出发绕回路一圈回来，电场对你做了净功，你没法给这一点贴上一个固定的高度标签。这正是本书反复强调“电势、电势能、电压不可混为一谈”的一个深层原因：电势这个概念只属于有头有尾的静电场。

至于方向，回路中感应电流的方向由一条叫\textbf{楞次定律}的规则给出：感应电流的效果，总是去\emph{阻碍}引起它的那个磁通量变化。磁铁靠近，回路就推开它；磁铁远离，回路就挽留它。这个“对着干”的脾气不是回路的任性，背后站着能量守恒，我们在数学翻译器里把这笔账算清楚。

\step{一句话模型}

\begin{oneline}
穿过一个回路的磁通量发生变化时，回路中就会出现感应电动势；从场的角度看，是变化的磁场在空间制造了一圈一圈的涡旋电场。感应电流的方向总是阻碍磁通量的变化——这正是能量守恒的要求。
\end{oneline}

\step{数学翻译器}

现在把上面的直觉一条一条翻译成可以计算的式子。

\textbf{第一条：磁通量。}面积为 \(S\) 的平面回路，放在磁感应强度为 \(B\) 的匀强磁场里，回路平面与磁场方向的夹角用回路的“正对程度”来描述：记回路法线（垂直于回路平面的方向）与磁场方向的夹角为 \(\theta\)，则磁通量

\[
\Phi = B\,S\cos\theta
\]

单位是韦伯（Wb）。式子左边回答“穿过回路的磁场总账有多少”；右边三项各自对应直觉里的三个因素：磁场强弱、回路大小、回路朝向。

\begin{center}
\begin{tikzpicture}[scale=0.9]
  % 匀强磁场，竖直向上
  \foreach \x in {0,0.8,1.6,2.4,3.2,4.0} {
    \draw[->] (\x,0) -- (\x,3.2);
  }
  \node[right] at (4.05,2.6) {$\mathbf B$};
  % 倾斜放置的回路平面
  \draw[thick,fill=gray!12] (0.8,0.7) -- (3.2,0.7) -- (3.9,1.9) -- (1.5,1.9) -- cycle;
  % 法线
  \draw[->,thick] (2.35,1.3) -- (2.05,3.0) node[above] {法线};
  % 角度标记：法线与磁场方向的夹角
  \draw (2.35,2.5) arc[start angle=90,end angle=100,radius=1.2];
  \node at (2.28,2.42) {$\theta$};
  \node[below] at (2.35,0.55) {回路面积 $S$};
\end{tikzpicture}
\end{center}

公式一写出来，立刻用特殊情况检查它。若回路正对磁场（\(\theta=0\)），\(\cos\theta=1\)，通量最大，等于 \(BS\)——渔网正对着水流，兜住最多。若回路平躺在磁场里（\(\theta=90^{\circ}\)），\(\cos\theta=0\)，通量为零——网面顺着水流，什么都兜不住。把 \(B\) 加倍或者把面积加倍，通量都正好加倍，和直觉一致。合理。

\textbf{第二条：法拉第电磁感应定律。}磁通量变化得越快，感应电动势越大；回路绕了 \(N\) 匝，每一匝都被推一把，效果乘以 \(N\)。写出大小关系：

\[
\mathcal{E} = N\,\frac{\Delta\Phi}{\Delta t}
\]

左边是感应电动势，单位仍是伏特；右边是每匝回路磁通量的变化率。注意式子里出现的是通量的\emph{变化}除以\emph{时间}，而不是通量本身——这是本章最容易记错的一点，我们马上会专门讨论。

照例检查特殊情况。磁铁停在线圈里不动：\(\Delta\Phi=0\)，于是 \(\mathcal{E}=0\)——再强的磁铁、再大的通量，不变就没有感应。把磁铁\emph{反向}插一次，\(\Delta\Phi\) 变号，电动势方向反转，电流跟着反向——这正是软铁环实验里“接通摆一下、断开反向摆一下”的指针。变化的时间 \(\Delta t\) 越短，电动势越大：同样一块磁铁，猛插一下得到的电压是慢慢插入的好几倍。方向、零值、快慢，全部对得上。

\begin{mathbridge}
用导数的语言，法拉第定律写成更紧凑的形式：

\[
\mathcal{E} = -N\frac{\mathrm{d}\Phi}{\mathrm{d}t}
\]

负号不是装饰，它就是楞次定律的数学化身：感应电动势的方向永远与磁通量的变化“对着干”。判断方向时与其背口诀，不如每次都问自己：感应电流产生的磁场，是在补正在减小的通量，还是在顶正在增大的通量？

还有一种常见情形可以直接推出结果：一段长 \(L\) 的直导线，以速度 \(v\) 垂直切割匀强磁场 \(B\)。导线里的自由电荷跟着导线一起运动，洛伦兹力沿导线方向的分量把正电荷推向一端、负电荷推向另一端，直到两端积累的电荷产生的电场力与磁力平衡。平衡条件 \(qE=qvB\) 给出两端电压

\[
\mathcal{E} = B\,L\,v
\]

检查：\(v=0\) 时为零，合理；\(B\) 反向则电动势反向，合理；导线越长、磁场越强、运动越快，电压越大，合理。有趣的是，这个结果和“回路面积变化导致通量变化”的算法给出完全相同的答案——同一个事实的两种记账方式。这个巧合其实藏着深意，我们在“模型的边界”里再回来谈。

如果让线圈在匀强磁场里以角速度 \(\omega\) 匀速转动，通量就是 \(\Phi=BS\cos\omega t\)，对它求导即得

\[
\mathcal{E} = NBS\omega\sin\omega t
\]

电动势随时间按正弦规律变化——这就是交流电的诞生方式。下面那幅图画的正是它。
\end{mathbridge}

\begin{center}
\begin{tikzpicture}[scale=1.0]
  % 坐标轴
  \draw[->] (0,0) -- (7.2,0) node[right] {$t$};
  \draw[->] (0,-1.6) -- (0,1.8) node[above] {$\mathcal{E}$};
  % 正弦曲线：电动势随时间正负交替
  \draw[thick,domain=0:6.8,samples=80] plot (\x,{1.3*sin(\x r)});
  % 标注峰值与零点
  \draw[dashed] (1.5708,0) -- (1.5708,1.3);
  \node[above] at (1.5708,1.3) {线圈与磁场平行：通量为零，变化最快};
  \node[below right] at (3.2,-0.15) {线圈正对磁场：通量最大，变化率为零};
  \draw[dashed] (3.1416,-0.6) -- (3.1416,0);
\end{tikzpicture}
\end{center}

这幅图值得停下来看懂，因为它纠正一个极其普遍的直觉错误。曲线过零的时刻，正是线圈正对磁场、磁通量\emph{最大}的时刻——但此时通量曲线正处在“山顶”，变化率为零，所以电动势为零。反过来，电动势达到峰值的时刻，是线圈平躺、磁通量为\emph{零}的时刻——此时通量变化最猛烈。\textbf{感应电动势看的是通量的变化率，不是通量本身。}磁场再强、通量再大，不变就没有感应；通量再小，变得快就有强感应。这就是为什么“强磁场中匀速平移的线圈毫无电流”——它是本章送给直觉的一记反例：直觉以为“磁场强就该有电”，规律说“变才有，不变就没有”。

\textbf{第三条：楞次定律与能量账。}现在把“阻碍变化”这句话换成能量语言。设想楞次定律反过来：感应电流不是阻碍而是\emph{帮助}磁通变化。那么把一块磁铁轻轻推向一个闭合线圈，感应电流会把它往里吸，磁铁加速，电流增大，吸力更强，磁铁更快……你什么都没有付出，磁铁却越跑越快、电流越来越大——一台零投入的永动机。能量守恒不允许这样的事发生，所以感应电流只能对着干：你想推进磁铁，它产生磁场推你的磁铁；你想抽出磁铁，它产生磁场挽留。你在手电筒手柄上感到的那股“发涩”的阻力，正是线圈里的感应电流通过磁场反过来咬住了你手里的磁铁。

这笔账的妙处在于它是闭合的：你克服阻力做的功，恰好等于回路里电流最终耗散掉的能量（加上储存起来的部分）。楞次定律不是一条新的力学定律，而是能量守恒在电磁感应里的“出纳员”。

\textbf{第四条：一个带真实数据的估算。}手摇发电手电筒的内部大致是：永磁体表面磁场约 \(B\approx 0.3\) 特斯拉（钕铁硼磁铁的典型值），线圈面积约 \(S\approx 2\ \mathrm{cm^2}=2\times10^{-4}\ \mathrm{m^2}\)，匝数约 \(N=500\)，手柄通过齿轮组把磁体增速到每秒约 20 转，即 \(\omega\approx 2\pi\times20\approx 126\) 弧度每秒。峰值电动势

\[
\mathcal{E}_{\max}=NBS\omega \approx 500\times0.3\times2\times10^{-4}\times126\approx 3.8\ \mathrm{V}
\]

点亮一颗白色发光二极管需要约 3 伏——数量级正好够。摇动快慢的变化通过 \(\omega\) 直接进入电压：摇得慢，电压掉到二极管门槛以下，灯就暗。这也回答了“先猜一猜”里的定量问题：摇快一倍，电压大约高一倍，但灯珠亮度并不简单翻倍，因为电流还受电路中其他元件的限制。

\textbf{第五条：线圈也会“自己感应自己”——自感。}到此为止磁通变化都来自外部，但请想一步：线圈通了电流，自己就会产生磁场，这个磁场同样穿过线圈自己。那么当电流变化时，线圈自身穿过的磁通也在变，于是线圈在自己身上感生出电动势，阻碍自己电流的变化。这叫\textbf{自感}。写成式子：自感电动势正比于电流的变化率，比例系数叫电感 \(L\)，单位是亨利：

\[
\mathcal{E}_{L} = L\,\frac{\Delta I}{\Delta t}
\]

电感在电路里的脾气很像力学里的惯性：质量大的物体讨厌速度突变，电感大的线圈讨厌电流突变。这个比喻像的地方在于：两者都储存能量（动能与磁场能），都让“状态量”只能连续变化、不能跳变，你要让物体跑起来或者让线圈里的电流建立起来，都得花一段时间、做一笔功。不像的地方在于：惯性质量储存的是物体运动的能量，而电感储存的是磁场的能量，线圈里的电子并没有因此获得什么“机械惯性”；一旦电流稳定下来，电感就完全安静，不再插手电路的事。

自感不是书斋里的概念。老式日光灯靠一个叫镇流器的大电感工作：启动瞬间，启辉器断开，电感中的电流被迫急变，感生出上千伏的高压，把灯管里的气体击穿点亮。家里开关断开时偶尔能看到的小火花，也是同一个道理——电流被强行快速切断，线路的自感感生出高电压，把空气击穿。

\textbf{第六条：两个线圈互相感应——互感与变压器。}把两个线圈绕在同一个铁芯上，一个线圈（原线圈）通上变化的电流，铁芯里的磁场跟着变，变化全部穿过另一个线圈（副线圈），副线圈就得到感应电动势。铁芯的作用是把磁场“关”在一条闭合通道里，让原线圈产生的磁通几乎一丝不漏地穿过副线圈。理想变压器的电压比等于匝数比：

\[
\frac{U_1}{U_2}=\frac{N_1}{N_2}
\]

检查特殊情况：副线圈匝数为零，输出为零，合理；匝数比一比一，电压不变，合理。你的手机充电器里就有一个小变压器，把 220 伏市电先降到几伏：按匝数比，原副线圈匝数大约是 \(44:1\) 才能从 220 伏得到 5 伏。小区电线杆上的变压器深夜发出的嗡嗡声，是铁芯里的硅钢片被每秒变化一百次的磁场反复磁化、轻微伸缩振动发出的——你听到的其实就是“变化的磁场”本身的声音。

\textbf{第七条：直流与交流，以及电路的下一座桥。}电池提供的是直流：电压方向固定，电荷沿一个方向持续流动。而旋转线圈发出的电，方向每秒来回颠倒——我国市电每秒完成 50 个周期，电流方向每秒反转 100 次，这叫交流。电网为什么用交流？核心原因就在本章：\emph{变压器只对变化的磁场工作}，直流产生恒定磁场，通进变压器原线圈，副线圈除了接通和断开那两瞬之外什么也得不到。只有交流能被变压器随意升压降压——发电厂升压到几十万伏远距离送电（电流小，导线上的焦耳热损耗小），到小区门口再降压。电磁感应不只是发电的原理，也是输电可行、经济的原因。

说起发电，值得停下来看一眼发电厂在干什么。除了太阳能光伏板，人类几乎所有的发电方式——水电站、火电站、核电站、风力机——最后一道工序都是同一件事：让磁铁和线圈持续地相对转动。水轮机被江水推着转，汽轮机被高压蒸汽推着转，风车被风推着转；火电站烧煤、核电站烧铀，烧出来的热量先去烧开水产生蒸汽，蒸汽再去推汽轮机。听起来曲折得可笑，但每一步都有不得不如此的理由：煤炭和铀提供的是热能，热能没法直接变成电流，而“转动的线圈加磁场”是目前人类掌握的规模上最划算的能量转换接口。一座百万千瓦的电站，本质上就是一台被江河或蒸汽驱动的、巨大的本章演示实验。

电路里有了电感这个新角色，三个“基本演员”就凑齐了：电阻像摩擦，把电能不可逆地耗散成热；电容像弹簧（第 21 章），储存电荷和电场能；电感像惯性飞轮，储存磁场能。它们两两组合就引出本书后面的桥梁内容：电阻加电容（RC）决定电容充电放电的快慢，是相机闪光灯“充一会儿才能闪”的原因；电阻加电感（RL）决定电流建立和消失的快慢；而电感加电容再加一点电阻（RLC），就是“弹簧加惯性加阻尼”的电学翻版——电流在电容和电感之间来回振荡，电场能和磁场能互相转化，如同秋千在动能和势能之间来回。收音机选台，选的正是这个振荡的回音。这条线索会在讨论麦克斯韦方程和电磁波（第 24、25 章）时再次拾起。

\step{模型的边界}

每一条重要公式都要回答“何时成立、何时失效”，电磁感应定律尤其值得认真回答，因为它的成立范围出人意料地广，而它的常见误用也出奇地多。

\textbf{成立范围}。作为实验事实，法拉第定律对宏观回路普遍成立：磁通变化可以由磁场强弱变化、回路面积变化、回路朝向变化、回路在磁场中移动中的任何一种或几种共同引起，定律一概适用，与材料无关、与温度无关、与回路里接不接负载无关。从地下观测站的地磁记录线圈到粒子加速器里精密的感应探头，它都在工作。它没有“磁场太强就失效”的上限——医院核磁共振仪里几特斯拉的强磁场、脉冲星周围难以想象的磁场，都服从同一条定律。

\textbf{误用一：把“有磁场”当成“有感应”。}前面已经强调过，但值得再立一块警示牌：恒定的强磁场中，静止或匀速平移的线圈里没有任何感应电流。判断感应的唯一问题是“穿过回路的磁通量在变吗”，而不是“磁场强不强”。这个错误直觉的根源，大概是我们把“强”和“活跃”混为一谈——一座静止的高山很雄伟，但它推不动任何东西。

\textbf{误用二：把楞次定律说成“阻碍磁铁运动”。}更准确的说法是阻碍磁通量的\emph{变化}。磁铁进入线圈时被排斥，离开时被吸引；但如果磁铁静止地待在线圈正中央，什么力也没有。楞次定律盯着的是“变”，不是“动”——只是多数情况下，动恰好意味着变。

\textbf{误用三：以为磁通定律能解决一切回路问题。}有一个著名的反例叫法拉第圆盘（它也是人类第一台发电机）：一个铜圆盘在磁场中转动，转轴和盘边各接一个电刷引出电流。盘的位置不变、穿过整个装置的磁通似乎也不变，可电流实实在在流出来了。问题出在“回路”的选取上：电流的路径有一部分在转动的盘内，随盘一起切割磁场。严格说，这时用“导线切割磁感线”（动生电动势）的算法才清楚，生搬“回路磁通变化”会把自己绕进矛盾。这不是定律错了，而是提醒我们：磁通定律里“回路”这个词，在含有运动导体边界的情形下需要小心定义。

\textbf{更深处：两种感应其实是同一件事。}回头看，磁通变化有两种来源：线圈在磁场中运动（动生，可用洛伦兹力解释）和磁场本身变化（感生，必须用涡旋电场解释）。但“谁在动”取决于你站在哪个参考系里看：你拿着磁铁冲向静止线圈，和你拿着线圈冲向静止磁铁，对线圈来说效果完全一样。爱因斯坦在 1905 年创立相对论的论文，开头讨论的正是这个“磁铁与导体”的不对称：同一件事，旧理论却给出两套不同的解释。他要求自然界只保留一套解释，结果逼出了狭义相对论——电场和磁场在不同参考系里会互相转化。这段故事我们在第 38 章继续。

\textbf{本章埋下的最大伏笔}。本章说：变化的磁场制造电场。麦克斯韦盯着这句话看了一会儿，问了一个对称性的问题：那变化的电场呢？他大胆补上另一半——变化的电场也制造磁场。这一补丁直接预言了电磁波的存在，并最终让人类认出：光就是电磁波。这是第 24、25 章的故事。

\begin{challenge}
把涡旋电场的性质写成场的语言：沿任意闭合回路，电场沿路径的累积（环流）等于穿过该回路的磁通量变化率的负值，

\[
\oint\mathbf{E}\cdot\mathrm{d}\mathbf{l}=-\frac{\mathrm{d}\Phi_B}{\mathrm{d}t}
\]

静电场里这个环流恒等于零，所以静电场是保守场，可以定义电势；涡旋电场环流不为零，所以它不是保守场，“电势”对它失效。这个对比是麦克斯韦方程组里“法拉第方程”的完整内容，也是理解第 24 章位移电流补丁的钥匙。另一个有趣的问题是：既然涡旋电场的电场线是闭合圆环，那么一块远离一切导线、真空中正在变化的磁场区域，周围真的存在电场吗？实验的回答是“真的存在”——电子感应加速器正是用环形真空室中变化的磁场把电子加速到接近光速，中间没有任何导线碰到电子。
\end{challenge}

\step{回头解释开场现象}

现在带着完整模型回到开头的三个现象，逐个结账。

\textbf{手摇手电筒}。你摇动手柄，齿轮组带着永磁体在线圈旁高速旋转，穿过线圈的磁通量以每秒几十次的频率反复增减换向，线圈中感生出方向反复变化的电动势——一小股交流电。电路里的二极管只允许一个方向的电流通过，把这股交流“捋直”，充进一颗小电容（所以摇完松开手灯还能亮几分钟：那是电容里存的电荷在继续供电）。而你手上感到的阻力，是楞次定律亲自登门：线圈里一旦有了电流，电流的磁场就对旋转的磁铁施加反抗的力矩，你克服它做的功，经由“机械能—电能—光能和热”的链条，一分不少地变成了灯的光。灯接得越亮，需要的电流越大，反抗力矩越大，摇起来越费力——能量守恒不接受讨价还价。

\textbf{磁铁穿过铝管}。磁铁下落时，管壁上每一小段圆环看到的磁通都在变化：磁铁上方，通量在减小；磁铁下方，通量在增大。铝是优良导体，变化的磁通在管壁里感生出一圈一圈的电流，它们有个形象的名字——\textbf{涡流}，像水流里的漩涡。按楞次定律，涡流的磁场永远阻碍磁铁的运动：在磁铁上方拉它一把，在下方顶它一下。磁铁受到向上的合力，只能缓缓飘落。它少掉的重力势能去了哪里？变成了涡流在铝管电阻上产生的焦耳热——铝管摸起来会微微发热，只是热量太小不易察觉。铜铝不吸磁铁，却能让磁铁“刹车”：它们互动的桥梁不是磁性，而是电磁感应。开头“先猜一猜”里选“一样快”的直觉，错的根源是把“没有磁吸引”误当成了“没有相互作用”。

\textbf{电磁炉}。玻璃面板下面藏着一只扁平的铜线圈，通着每秒变化两万次以上的高频交流。线圈上方如果放着铁锅，锅底金属中感生出强劲的涡流，电流遇到电阻，直接在锅底内部发热——热是在锅里生的，玻璃面板几乎不导电，自身不产生涡流，所以只是被锅“烤”温的。这就是为什么锅里的水开了，面板却不太烫。纯铜锅和铝锅电阻率太低，涡流虽大发热效率却差，多数电磁炉检测到这类锅具会干脆拒绝工作。同一个原理，反过来用就是金属探测器：安检门的线圈发出变化的磁场，你口袋里的钥匙中感生涡流，涡流的磁场反过来被仪器捕捉到——安检门“听见”的，是你身上感应电流的回答。

\textbf{无线充电与电吉他}。开场还提到手机放在充电板上“隔空”充电。现在你能看穿它了：充电板里是一只线圈，通着高频交流；手机背面贴着另一只线圈。两只线圈隔着几毫米的空气构成一个“没有铁芯的变压器”——互感。正因为没有铁芯约束磁场，漏掉的磁通很多，所以手机必须贴紧、对准才能高效充电，拿远一点效率就断崖式下跌。变压器的原理，在掌心大小的地方又演了一遍。还有一件乐器也值得回头看一眼：电吉他的拾音器，就是一块小磁铁外面绕着几千匝细线圈，固定在琴弦下方。琴弦是铁磁材料做的，它一振动，磁铁附近的磁场分布就跟着变化，穿过线圈的磁通随之起伏，线圈两端便送出与振动同步的电信号——琴弦的机械振动先变成磁通的变化，再变成电流的起伏，最后经放大器变成声音。吉他这位老朋友（它在本书讲张力和驻波时已经登场过）在电磁学里又露出了新的一层。

\step{脑洞实验与挑战题}

\begin{thought}[太空里的发电缆绳]
把电磁感应推到极端尺度：一根二十千米长的导线，从绕地球飞行的航天器上垂下去，会怎样？导线随航天器以每秒约 7.7 千米的速度横扫地球磁场（约 \(3\times10^{-5}\) 特斯拉），按动生电动势估算：

\[
\mathcal{E}=BLv\approx 3\times10^{-5}\times2\times10^{4}\times7.7\times10^{3}\approx 4.6\times10^{3}\ \mathrm{V}
\]

几千伏的电压，足够击穿空气了！这件事不是纯粹的脑洞：1996 年，美国航天飞机真的做过“绳系卫星”实验，用一根约二十千米长的缆绳拖着一颗卫星，计划在太空中发电。实验在缆绳放出大部分后测到了几千伏量级的电压，可惜缆绳随后意外断裂，卫星飘走。请思考：就算缆绳不断，电流从哪里来、流到哪里去？导线两端之间并没有连成回路——电流的回路要靠地球高层大气里稀薄的等离子体来闭合：缆绳一端从等离子体里收集电子，另一端通过电子枪把电子射回太空。你看，哪怕在四百千米高的轨道上，“电流必须成环”依然是铁律。
\end{thought}

\begin{puzzles}
\item 一个 \(N\) 匝、电阻为 \(R\) 的闭合线圈，穿过它的磁通量从 \(\Phi_1\) 变到 \(\Phi_2\)。第一次用一秒慢慢完成，第二次用百分之一秒猛地完成。两次过程中流过线圈某一截面的\emph{总电荷量}，哪次多？（提示：电流 \(I=\mathcal{E}/R\)，而总电荷是电流对时间的累积。把 \(\mathcal{E}=N\Delta\Phi/\Delta t\) 代进去，看看 \(\Delta t\) 会不会约掉。快拉电流大但时间短，慢拉电流小但时间长——算算总账。）
\item 两根完全一样的竖直铝管，其中一根沿侧面从头到尾开了一条细细的竖缝。两块完全一样的磁铁同时从管口放入，哪根管里的磁铁先落地？（提示：涡流需要在管壁上绕成\emph{闭合}的圆圈。一条竖缝对圆圈意味着什么？再想一步：开缝管里的磁铁几乎自由落体，那“先猜一猜”里你的直觉错在哪一层？）
\item 变压器的铁芯为什么不用一整块铁，而是用许多涂了绝缘漆的薄硅钢片叠起来？（提示：铁芯里的磁场每秒变化一百次，整块铁内部本身就是导体。顺着磁场方向看进去，整块铁里能容纳多大的涡流回路？切成薄片、片间绝缘之后，回路被切成了什么样？这和挑战 2 里的竖缝是同一个思想。）
\item 假设地球磁场（约 \(5\times10^{-5}\) 特斯拉）在一分钟内整体翻转方向。你书桌上一个 100 匝、面积 1 平方米的废弃线圈两端会出现多大电压？会电死人吗？（提示：翻转意味着通量从 \(+BS\) 变到 \(-BS\)，\(\Delta\Phi\) 是 \(2BS\)。算出数量级再下结论——电磁感应无处不在，但“有没有效应”和“效应有多大”是两个独立的问题。）
\end{puzzles}
